% Source: https://artofproblemsolving.com/community/c3h1533911_minimum_number_of_teachers
% Stable identity: aops-topic-1533911-q01
\textbf{Lower bound.} In each period all $463$ students are in classes of at most $28$, so there are at least $\left\lceil\dfrac{463}{28}\right\rceil=17$ classes in each period ($16\cdot28=448<463$). That is at least $6\cdot17=102$ classes in all, and with at most $5$ classes each, at least $\left\lceil\dfrac{102}{5}\right\rceil=21$ instructors are needed.

\textbf{$21$ instructors suffice.} Run exactly $17$ classes in each period (for example $9$ classes of $27$ and $8$ of $28$: $243+224=467\ge463$, so the students fit). List the $102$ classes period by period and assign instructors $1,2,\dots,21,1,2,\dots,21,\dots$ cyclically. The $17$ classes of one period are consecutive in the list, so they get $17$ different instructors (as $17\le21$), and each instructor gets at most $\left\lceil\dfrac{102}{21}\right\rceil=5$ classes.

So the minimum is $\boxed{21}$.
